\section[Descenso de propiedades conjuntistas]{Descenso de propiedades
conjuntistas y de propiedades de finitud de morfismos\protect\footnotemark}
\label{VIII.3}


\footnotetext{Para otros resultados como los tratados en los números 3
y~4, \Cf EGA~IV~2.3, 2.6 y~2.7.}

\begin{proposition}
\label{VIII.3.1}
Sean $f\colon X\to Y$ un~$S$-morfismo, $g\colon S'\to S$ un morfismo
sobreyectivo y $f'\colon X'=X\times_{S}S'\to Y'=Y\times_{S}S'$ el morfismo
deducido de~$f$ mediante el cambio de base por~$g\colon S'\to S$. Para que
$f$ sea sobreyectivo (\resp radicial), es necesario y suficiente que lo
sea~$f'$.
\end{proposition}

Obsérvese que $f'$ también puede obtenerse mediante el cambio de base
$Y'\to Y$, que asimismo es sobreyectivo, pues se deduce del morfismo
sobreyectivo~$g\colon S'\to S$. Por otra parte, para todo~$y\in Y$ y
todo~$y'\in Y'$ situado sobre~$y$, se tiene un isomorfismo
$$
X'_{y'}\isomfrom X_{y}\otimes_{\kres(y)}\kres(y')\quoi,
$$
donde $X_{y}$ denota la fibra de~$X$ en~$y$, y $X'_{y'}$ la de~$X'$ en
$y'$. Se sigue que~$X_{y}$ es no vacía (\resp tiene a lo sumo un punto, y
este corresponde a una extensión residual radicial) si y solo si
\ifthenelse{\boolean{orig}}
{$X_{y'}$}
{$X'_{y'}$}
tiene la misma propiedad. Esto demuestra~\Ref{VIII.3.1}.

\begin{corollaire}
\label{VIII.3.2}
Bajo las condiciones de~\Ref{VIII.3.1}, si $f'$ es inyectivo
(\resp biyectivo), entonces $f$ también lo es.
\end{corollaire}

Esto se debe a que, si $X'_{y'}$ tiene a lo sumo un punto
(\resp exactamente un punto), lo mismo ocurre con~$X_{y}$; así es, en
efecto, puesto que el morfismo $X'_{y'}\to X_{y}$ es sobreyectivo (por
deducirse del morfismo sobreyectivo
$\Spec(\kres(y'))\to \Spec(\kres(y))$).

\begin{proposition}
\label{VIII.3.3}
Con las notaciones de~\Ref{VIII.3.1}, supongamos que~$g\colon S'\to S$ es
sobreyectivo y cuasicompacto (\resp fielmente plano y cuasicompacto). Para
que~$f$ sea cuasicompacto (\resp de tipo finito), es necesario y suficiente
que lo sea~$f'$.
\end{proposition}

\marginpar{204}
Solo hay que demostrar la suficiencia. Puede suponerse, evidentemente,
$S=Y$, pues la hipótesis hecha sobre~$g\colon S'\to S$ se conserva para
$Y'\to Y$. Además, puede suponerse que~$Y$ es afín. Entonces $Y'$ es
cuasicompacto y, por tanto, $X'$ es cuasicompacto (pues~$f'$ lo es por
hipótesis). Sea $(X_{i})_{i\in I}$ una familia de abiertos afines de~$X$ que
recubre~$X$; los~$X_{i}'$ son abiertos de~$X'$ que recubren~$X'$, por lo que
una subfamilia finita recubre~$X'$. Como $X'\to X$ es sobreyectivo, se sigue
que los~$X_{i}$ correspondientes ya recubren~$X$; por tanto, $X$ es
cuasicompacto, es decir, $f$ es cuasicompacto. Supongamos ahora que~$f'$ es
de tipo finito y demostremos que~$f$ lo es, suponiendo~$g$ fielmente plano.
Reemplazando $Y'$ por la suma de una familia de abiertos afines que lo
recubre, puede suponerse que~$Y'$ es afín. Por último, como, por lo que
precede, $X$ está recubierto por un número finito de abiertos afines
$X_{i}$, hay que demostrar que estos son de tipo finito sobre~$Y$, sabiendo
que $X_{i}'$ es de tipo finito sobre~$Y'$. Quedamos reducidos entonces al

\begin{corollaire}
\label{VIII.3.4}
Sean $B$ una $A$-álgebra, $A'$ una~$A$-álgebra fielmente plana y
$B'=B\otimes_{A}A'$ la
\ifthenelse{\boolean{orig}}
{$A$-álgebra}
{$A'$-álgebra}
deducida de~$B$ por cambio de anillos. Para que~$B$ sea de tipo finito, es
necesario y suficiente que lo sea~$B'$.
\end{corollaire}

Solo hay que demostrar la suficiencia. Se tiene
$B=\varinjlim_i B_{i}$, donde los~$B_{i}$ son las subálgebras de tipo finito
de~$B$. Por tanto, $B'=\varinjlim_i B_{i}'$ y, si $B'$ es de tipo finito
sobre~$A'$, entonces $B'$ es igual a uno de los~$B_{i}'$; por consiguiente,
$B$ es igual a~$B_{i}$ y, por tanto, es de tipo finito.

\begin{corollaire}
\label{VIII.3.5}
Supongamos de nuevo que el morfismo de cambio de base~$g\colon S'\to S$ es
fielmente plano y cuasicompacto. Para que~$f$ sea cuasifinito, es necesario y
suficiente que lo sea~$f'$.
\end{corollaire}

En efecto, la propiedad «ser cuasifinito» es, por definición, la conjunción
de «ser de tipo finito» y «tener fibras finitas», cada una de las cuales
desciende bien por~$g$: la primera en virtud de~\Ref{VIII.3.3}, y la segunda
por el razonamiento de~\Ref{VIII.3.1} (que solo utiliza la sobreyectividad
de~$g$).

\begin{remarques}
\label{VIII.3.6}
Sean $A$ un anillo y $X$ un~$A$-preesquema. Se ve fácilmente que
\marginpar{205}
las condiciones siguientes son equivalentes:

\begin{enumerate}
\item[(i)] Existe un anillo noetheriano~$A_{0}$ (que, si se desea, puede
suponerse un subanillo de tipo finito de~$A$), un $A_{0}$-preesquema de tipo
finito~$X_{0}$, un homomorfismo
\ifthenelse{\boolean{orig}}
{$A\to A_{0}$}
{$A_0\to A$}
y un $A$-isomorfismo~$X\isomto X_{0}\times_{A_{0}}A$.
\item[(ii)] El morfismo diagonal~$X\to X\times_{\Spec(A)}X$ es cuasicompacto
(condición vacía si $X$ está separado sobre~$A$), y $X$ es reunión finita de
abiertos afines~$X_{i}$ cuyos anillos~$B_{i}$ son álgebras de presentación
finita sobre~$A$, es decir, cocientes de álgebras de polinomios en un número
finito de indeterminadas por ideales de tipo finito.
\end{enumerate}

Si $X$ es él mismo afín, de anillo~$B$, estas condiciones significan
simplemente que~$B$ es un álgebra de presentación finita sobre~$A$.

Un morfismo~$f\colon X\to Y$ se llama \emph{morfismo de presentación
finita},
\index{presentacion finita@presentación finita (morfismo de)|hyperpage}%
y se dice también que~$X$ es de presentación finita sobre~$Y$ si $Y$ es
reunión de abiertos afines~$Y_{i}$
\ifthenelse{\boolean{orig}}
{tal que}
{tales que}
$X|Y_{i}$, como $Y_{i}$-preesquema, satisface las condiciones equivalentes
anteriores. Lo mismo ocurre entonces con $X|Y'$ para \emph{todo} abierto
afín~$Y'$ de~$Y$. Esta propiedad es estable por cambio de base y, además, el
compuesto de dos morfismos de presentación finita es de presentación finita.

Sentadas estas nociones, se ve en~(ii), procediendo como
en~\Ref{VIII.1.10}, que este enunciado sigue siendo válido al reemplazar en
él las palabras «de tipo finito» por «de presentación finita».
\end{remarques}
